\section{Additional Spectral Comparisons}
\label{sec:additional-spectral-comparisons}
\label{sec:fixed-amplitude-spectra}

\subsection{Fixed-amplitude spectra}

Figure~\ref{fig:simulation-results} resolves the narrow-detuning maps of
main-text Fig.~2 into five measured amplitude slices per protocol. The
nearest acquired rows to nominal peak Rabi frequencies of 3, 10, 20, 30,
and $40~\mathrm{MHz}$ are selected without optimizing their agreement with
the model. The three-level calculation uses the exact acquired drive values
and all 501 measured detunings over $\pm0.5~\mathrm{MHz}$, with the same
primary-qubit model parameters as main-text Fig.~2. Each measured point
contains 2000 repetitions. Halving the integration step at all five drives
and seven representative detunings changes $P_e$ by less than
$1.2\times10^{-5}$. This checks numerical accuracy independently of
the model's agreement with experiment.

To compare line shapes independently of their absolute offsets, we subtract
each trace's mean over the full displayed detuning interval and add the same
stacking offset to its experimental and simulated members. The vertical scale
is common to every trace: no contrast normalization, gain adjustment,
frequency shift, or smoothing of the measured data is applied. This display
retains differences in feature depth, width, and center, while removing only
a constant population background. The unaligned maps remain in the main text.

\begin{figure*}[!tp]
    \centering
    \includegraphics[width=\textwidth]{%
        ../figures/paper/supplemental/s02_spectrum_slices.pdf}
    \caption{Stacked experimental and simulated spectra for
    (a) root-Lorentzian and (b) echo-root-Lorentzian pulses, with
    $L=20~\mu\mathrm{s}$ and $c=0.005$. From bottom to top, the acquired
    peak Rabi frequencies are approximately 2.92, 9.95, 20.07, 30.03,
    and $39.98~\mathrm{MHz}$, labeled individually. Colored markers show
    all measured points from the narrow-detuning scans in main-text Fig.~2;
    black lines show three-level simulations at the identical drive and
    detuning values. For each trace independently,
    $y=P-\langle P\rangle_{\Delta}+0.55j$, where $j=0,\ldots,4$
    indexes increasing drive. The scale bars denote a population difference
    of 0.2. Only constant vertical offsets are removed; contrast and frequency
    are not fitted. Two-state readout does not independently resolve
    $|f\rangle$, whereas the simulated observable is $P_e$.
    }
    \label{fig:simulation-results}
\end{figure*}


\subsection{Duration--contrast tradeoff}

Figure~\ref{fig:duration-resolution-comparison} quantifies the tradeoff between
pulse duration, linewidth, and signal contrast at fixed cutoff $c=0.005$.
The continuous curves show the simulated response, while the markers show the
experimental data.  Panels (a) and (b) separate spectral narrowing from signal
loss: longer pulses reach smaller FWHM relative to the $T_2$ limit, but their
fitted contrast remains lower over the usable drive range.  Their
contrast-weighted operating windows in panel (c) therefore contract despite
the high resolution.  Pulse duration must consequently be selected jointly
with contrast and experimental signal quality, rather than by linewidth alone.

\begin{figure*}[!tp]
    \centering
    \includegraphics[width=\textwidth]{%
        ../figures/paper/supplemental/s03_duration_resolution.pdf}
    \caption{Resolution--contrast performance versus peak Rabi frequency up
    to $25~\mathrm{MHz}$ for echo-root-Lorentzian pulses with $c=0.005$ and
    $L=5$, 10, 15, 20, and $30~\mu\mathrm{s}$.  Colored continuous lines show
    two-level simulations, and colored markers show OPX1000 data.  (a) Fitted
    FWHM in units of the coherence-limited reference,
    $\Gamma_f/\Gamma_{f,T_2}$, with the dashed line marking
    $\Gamma_f/\Gamma_{f,T_2}=1$; (b) fitted contrast $A$; and (c)
    contrast-weighted resolution
    $Q=AR=A\Gamma_{f,T_2}/\Gamma_f$.  The red curve is the constant-drive
    reference from the Bloch equations, and the dashed line in (c) marks $Q=0.10$.}
    \label{fig:duration-resolution-comparison}
\end{figure*}

